\input{../tex-inputs/latex-korrekturansicht-vorspann}

\section[Theodor von Sosnosky an Arthur Schnitzler, 21. 6. 1901]{L04232 Theodor von Sosnosky an Arthur Schnitzler, 21. 6. 1901}
\nopagebreak\mylabel{L04232v}
\rehead{ }\normalsize\beginnumbering\briefempfaengerindex{Schnitzler, Arthur@\textsc{Schnitzler, Arthur}!zzzSosnosky, Theodor von@\emph{von Theodor von Sosnosky}!1901-06-212@{21. 6. 1901}|(be}
\toendnotes[C]{\smallbreak\pagebreak[2]}
\correspDesc{Versand  durch Theodor von Sosnosky am 21. 6. 1901 in Alland
\newline{}Erhalt  durch Arthur Schnitzler im Zeitraum [22. 6. 1901
                  – 26. 6. 1901?] in Wien}\toendnotes[C]{\smallbreak}
\Standort{DLA, A:Schnitzler, HS.1985.1.4640.}
\physDesc{Brief, 1 Blatt, 2 Seiten, 3004 Zeichen
\newline{}Schreibmaschine
\newline{}Handschrift: blaue Tinte, lateinische Kurrent (\noindent{}Datumszeile, Unterschrift, Korrekturen, drei Unterstreichungen)
\newline{}Schnitzler: mit rotem Buntstift eine Unterstreichung }\toendnotes[C]{\smallbreak}
\pstart
           \raggedleft{}{\pb}\textcolor{pink}{Alland, N-Oest}\oindex{Alland@\textbf{Alland}, \emph{Verwaltungsgebiet}|pw}{}\ledrightnote{\textcolor{pink}{Alland}}{ }21./6. 901.\pend
           
\pstart\center{}{[}ms.:{]} Sehr geehrter Herr Doktor!\pend\vspace{0.5em}
\pstart
           Anbei erlaube ich mir, Ihnen die vorgestern im \textcolor{green}{N. W. T.}\pwindex{Neues Wiener Tagblatt@\emph{Neues Wiener Tagblatt}|pw}{}\ledrightnote{\textcolor{green}{Neues Wiener Tagblatt}} erschienene \label{K_L04232-1v}\edtext{\textcolor{green}{Besprechung}\pwindex{Sosnosky, Theodor von 4.\,1.\,1866 Budapest – 3.\,2.\,1943 Wien@\textsc{Sosnosky, Theodor von} (4.\,1.\,1866 Budapest – 3.\,2.\,1943 Wien), \emph{Schriftsteller}!Neue Erzählungsliteratur@\strich\emph{Neue Erzählungsliteratur}|pwv}{}\ledrightnote{{$\rightarrow$}\emph{\textcolor{green}{Neue Erzählungsliteratur}}}}{\lemma{\textnormal{\emph{Besprechung}}}\Cendnote{\textnormal{\textcolor{blue}{Theodor von Sosnosky}\pwindex{Sosnosky, Theodor von 4.\,1.\,1866 Budapest – 3.\,2.\,1943 Wien@\textsc{Sosnosky, Theodor von} (4.\,1.\,1866 Budapest – 3.\,2.\,1943 Wien), \emph{Schriftsteller}|pwk}: \emph{\textcolor{green}{Neue Erzählungsliteratur}\pwindex{Sosnosky, Theodor von 4.\,1.\,1866 Budapest – 3.\,2.\,1943 Wien@\textsc{Sosnosky, Theodor von} (4.\,1.\,1866 Budapest – 3.\,2.\,1943 Wien), \emph{Schriftsteller}!Neue Erzählungsliteratur@\strich\emph{Neue Erzählungsliteratur}|pwk}}. In: \emph{\textcolor{green}{Neues Wiener Tagblatt}\pwindex{Neues Wiener Tagblatt@\emph{Neues Wiener Tagblatt}|pwk}}, Jg. 35, Nr. 167,
                        20. 1. 1901, S. [1]–3. Es handelt sich um eine
                  Sammelrezension, die außer \emph{\textcolor{green}{Bertha Garlan}\pwindex{Schnitzler, Arthur 15. 5. 1862 Wien – 21. 10. 1931 ebd.@\textsc{Schnitzler, Arthur} (15. 5. 1862 Wien – 21. 10. 1931 ebd.), \emph{Schriftsteller, Mediziner}!Frau Bertha Garlan. Roman@\strich\emph{Frau Bertha Garlan. Roman}|pwk}}
                  verschiedene Neuerscheinungen von \textcolor{blue}{Georg von
                     Ompteda}\pwindex{Ompteda, Georg von 29.\,3.\,1863 Hannover – 10.\,12.\,1931 München@\textsc{Ompteda, Georg von} (29.\,3.\,1863 Hannover – 10.\,12.\,1931 München), \emph{Schriftsteller}|pwk}, \textcolor{blue}{Maximilian von Rosenberg}\pwindex{Rosenberg, Maximilian von 1849 Halberstadt – 1913@\textsc{Rosenberg, Maximilian von} (1849 Halberstadt – 1913)|pwk},
                     \textcolor{blue}{Wolf von Tainach}\pwindex{Tainach, Wolf von @\textsc{Tainach, Wolf von}, \emph{Schriftsteller}|pwk}, \textcolor{blue}{Emil Marriot}\pwindex{Marriot, Emil 20.\,11.\,1855 Wien – 1938 ebd.@\textsc{Marriot, Emil} (20.\,11.\,1855 Wien – 1938 ebd.), \emph{Schriftstellerin}|pwk} und \textcolor{blue}{Paul von
                     Schönthan}\pwindex{Schönthan-Pernwald, Paul von 19.\,3.\,1853 Wien – 4.\,8.\,1905 ebd.@\textsc{Schönthan-Pernwald, Paul von} (19.\,3.\,1853 Wien – 4.\,8.\,1905 ebd.), \emph{Schriftsteller, Journalist}|pwk} behandelt.}}}\label{K_L04232-1} von \textcolor{green}{Fr. B.
                  Garlan}\pwindex{Schnitzler, Arthur 15. 5. 1862 Wien – 21. 10. 1931 ebd.@\textsc{Schnitzler, Arthur} (15. 5. 1862 Wien – 21. 10. 1931 ebd.), \emph{Schriftsteller, Mediziner}!Frau Bertha Garlan. Roman@\strich\emph{Frau Bertha Garlan. Roman}|pw}{}\ledrightnote{\textcolor{green}{Frau Bertha Garlan. Roman}} zuübersenden. Ich muss jedoch bemerken, dass ich ausser und vor dieser
               noch eine zweite viel \textcolor{green}{ausführlichere}\pwindex{Sosnosky, Theodor von 4.\,1.\,1866 Budapest – 3.\,2.\,1943 Wien@\textsc{Sosnosky, Theodor von} (4.\,1.\,1866 Budapest – 3.\,2.\,1943 Wien), \emph{Schriftsteller}!Bücher und Zeitschriftenschau [Frau Bertha Garlan]@\strich\emph{Bücher und Zeitschriftenschau [Frau Bertha Garlan]}|pwv}{}\ledrightnote{{$\rightarrow$}\emph{\textcolor{green}{Bücher und Zeitschriftenschau [Frau Bertha Garlan]}}} geschrieben habe; mit der hatte ich aber arges Pech, zuerst
               nahm man sie nicht, weil sie zu freun\introOben{}d\introOben{}lich war, dann sandte
               ich sie der »\textcolor{green}{Nation}\pwindex{Nation. Wochenschrift für Politik, Volkswirtschaft und Literatur@\emph{Die Nation. Wochenschrift für Politik, Volkswirtschaft und Literatur}|pw}{}\ledrightnote{\textcolor{green}{Die Nation. Wochenschrift für Politik, Volkswirtschaft und Literatur}}« die das \textcolor{green}{Buch}\pwindex{Schnitzler, Arthur 15. 5. 1862 Wien – 21. 10. 1931 ebd.@\textsc{Schnitzler, Arthur} (15. 5. 1862 Wien – 21. 10. 1931 ebd.), \emph{Schriftsteller, Mediziner}!Frau Bertha Garlan. Roman@\strich\emph{Frau Bertha Garlan. Roman}|pwv}{}\ledrightnote{{$\rightarrow$}\emph{\textcolor{green}{Frau Bertha Garlan. Roman}}} aber leider schon vergeben hatte; so
               sandte ich die \textcolor{green}{Besprechung}\pwindex{Sosnosky, Theodor von 4.\,1.\,1866 Budapest – 3.\,2.\,1943 Wien@\textsc{Sosnosky, Theodor von} (4.\,1.\,1866 Budapest – 3.\,2.\,1943 Wien), \emph{Schriftsteller}!Bücher und Zeitschriftenschau [Frau Bertha Garlan]@\strich\emph{Bücher und Zeitschriftenschau [Frau Bertha Garlan]}|pwv}{}\ledrightnote{{$\rightarrow$}\emph{\textcolor{green}{Bücher und Zeitschriftenschau [Frau Bertha Garlan]}}} an
               die »\textcolor{green}{Norddeutsche Allgem.}\pwindex{Norddeutsche Allgemeine Zeitung@\emph{Norddeutsche Allgemeine Zeitung}|pw}{}\ledrightnote{\textcolor{green}{Norddeutsche Allgemeine Zeitung}}«, wo mir die Annahme
               ziemlich sicher schien; leider musste ich damit auf meinen Wunsch verzichten, den \textcolor{green}{Abdruck}\pwindex{Sosnosky, Theodor von 4.\,1.\,1866 Budapest – 3.\,2.\,1943 Wien@\textsc{Sosnosky, Theodor von} (4.\,1.\,1866 Budapest – 3.\,2.\,1943 Wien), \emph{Schriftsteller}!Bücher und Zeitschriftenschau [Frau Bertha Garlan]@\strich\emph{Bücher und Zeitschriftenschau [Frau Bertha Garlan]}|pwv}{}\ledrightnote{{$\rightarrow$}\emph{\textcolor{green}{Bücher und Zeitschriftenschau [Frau Bertha Garlan]}}} an auffallender Stelle
               gebracht zu wissen, denn dieses \textcolor{green}{Blatt}\pwindex{Norddeutsche Allgemeine Zeitung@\emph{Norddeutsche Allgemeine Zeitung}|pwv}{}\ledrightnote{{$\rightarrow$}\emph{\textcolor{green}{Norddeutsche Allgemeine Zeitung}}} pflegt Recensionen nicht im Feuilleton zubringen, sondern bloss in den
               Bücheranzeigen, das hat auch den Nachteil dass man nie erfährt, wann die \label{K_L04232-2v}\edtext{Arbeit erschienen}{\lemma{\textnormal{\emph{Arbeit erschienen}}}\Cendnote{\textnormal{Der Text erschien erst über ein halbes Jahr später unter
                  Kürzel und ohne definierenden Titel: \textcolor{blue}{Th. v. S.}\pwindex{Sosnosky, Theodor von 4.\,1.\,1866 Budapest – 3.\,2.\,1943 Wien@\textsc{Sosnosky, Theodor von} (4.\,1.\,1866 Budapest – 3.\,2.\,1943 Wien), \emph{Schriftsteller}|pwk}: \emph{\textcolor{green}{Bücher und Zeitschriftenschau}\pwindex{Sosnosky, Theodor von 4.\,1.\,1866 Budapest – 3.\,2.\,1943 Wien@\textsc{Sosnosky, Theodor von} (4.\,1.\,1866 Budapest – 3.\,2.\,1943 Wien), \emph{Schriftsteller}!Bücher und Zeitschriftenschau [Frau Bertha Garlan]@\strich\emph{Bücher und Zeitschriftenschau [Frau Bertha Garlan]}|pwk}}. In: \emph{\textcolor{green}{Norddeutsche Allgemeine Zeitung}\pwindex{Norddeutsche Allgemeine Zeitung@\emph{Norddeutsche Allgemeine Zeitung}|pwk}}, Nr. 50,
                        28. 2. 1902, Beilage.}}}\label{K_L04232-2} ist, denn das \textcolor{green}{Blatt}\pwindex{Norddeutsche Allgemeine Zeitung@\emph{Norddeutsche Allgemeine Zeitung}|pwv}{}\ledrightnote{{$\rightarrow$}\emph{\textcolor{green}{Norddeutsche Allgemeine Zeitung}}} schickt nur dann spontan
               Belege, wenn es sich um grössser\introOben{}e\introOben{} honorirte Artikel
               handelt. Da ich nun vor 8 Tagen erfuhr, dass das \textcolor{green}{N. W. T.}\pwindex{Neues Wiener Tagblatt@\emph{Neues Wiener Tagblatt}|pw}{}\ledrightnote{\textcolor{green}{Neues Wiener Tagblatt}} demnächst \label{T_L04232-1v}\edtext{eine}{\lemma{\textnormal{\emph{eine}}}\Cendnote{\textnormal{In der Vorlage steht:
                  »einen«.}}}\label{T_L04232-1}{ }\label{K_L04232-3v}\edtext{\textcolor{green}{Erzählung}\pwindex{Ompteda, Georg von 29.\,3.\,1863 Hannover – 10.\,12.\,1931 München@\textsc{Ompteda, Georg von} (29.\,3.\,1863 Hannover – 10.\,12.\,1931 München), \emph{Schriftsteller}!Lebensgefährtinnen@\strich\emph{Die Lebensgefährtinnen}|pwv}{}\ledrightnote{{$\rightarrow$}\emph{\textcolor{green}{Die Lebensgefährtinnen}}} von Br. \textcolor{blue}{Ompteda}\pwindex{Ompteda, Georg von 29.\,3.\,1863 Hannover – 10.\,12.\,1931 München@\textsc{Ompteda, Georg von} (29.\,3.\,1863 Hannover – 10.\,12.\,1931 München), \emph{Schriftsteller}|pw}{}\ledrightnote{\textcolor{blue}{Georg von Ompteda}}}{\lemma{\textnormal{\emph{Erzählung von Br. Ompteda}}}\Cendnote{\textnormal{Vom 23. 6. 1901 bis zum
                     27. 7. 1901 erschien im \emph{\textcolor{green}{Neuen
                     Wiener Tagblatt}\pwindex{Neues Wiener Tagblatt@\emph{Neues Wiener Tagblatt}|pwk}} der Roman \emph{\textcolor{green}{Die
                     Lebensgefährtinnen}\pwindex{Ompteda, Georg von 29.\,3.\,1863 Hannover – 10.\,12.\,1931 München@\textsc{Ompteda, Georg von} (29.\,3.\,1863 Hannover – 10.\,12.\,1931 München), \emph{Schriftsteller}!Lebensgefährtinnen@\strich\emph{Die Lebensgefährtinnen}|pwk}} von \textcolor{blue}{Georg von
                     Ompteda}\pwindex{Ompteda, Georg von 29.\,3.\,1863 Hannover – 10.\,12.\,1931 München@\textsc{Ompteda, Georg von} (29.\,3.\,1863 Hannover – 10.\,12.\,1931 München), \emph{Schriftsteller}|pwk} in 33 Folgen.}}}\label{K_L04232-3} bringen werde, \label{K_L04232-4v}\edtext{dessen jüngstes \textcolor{green}{Buch}\pwindex{Ompteda, Georg von 29.\,3.\,1863 Hannover – 10.\,12.\,1931 München@\textsc{Ompteda, Georg von} (29.\,3.\,1863 Hannover – 10.\,12.\,1931 München), \emph{Schriftsteller}!Monte Carlo@\strich\emph{Monte Carlo}|pwv}{}\ledrightnote{{$\rightarrow$}\emph{\textcolor{green}{Monte Carlo}}}}{\lemma{\textnormal{\emph{dessen jüngstes Buch}}}\Cendnote{\textnormal{\textcolor{blue}{Georg von Ompteda}\pwindex{Ompteda, Georg von 29.\,3.\,1863 Hannover – 10.\,12.\,1931 München@\textsc{Ompteda, Georg von} (29.\,3.\,1863 Hannover – 10.\,12.\,1931 München), \emph{Schriftsteller}|pwk}: \emph{\textcolor{green}{Monte Carlo. Roman}\pwindex{Ompteda, Georg von 29.\,3.\,1863 Hannover – 10.\,12.\,1931 München@\textsc{Ompteda, Georg von} (29.\,3.\,1863 Hannover – 10.\,12.\,1931 München), \emph{Schriftsteller}!Monte Carlo@\strich\emph{Monte Carlo}|pwk}}, Berlin: \emph{\textcolor{brown}{F. Fontane}\orgindex{F. Fontane@F. Fontane|pwk}}{ }1901.}}}\label{K_L04232-4} mir der \textcolor{brown}{Verleger}\orgindex{F. Fontane@F. Fontane|pw}{}\ledrightnote{\textcolor{brown}{F. Fontane}} kürzlich
               gesandt hatte, so benutzte ich die gute Gelegenheit, die \textcolor{green}{Besprechung}\pwindex{Sosnosky, Theodor von 4.\,1.\,1866 Budapest – 3.\,2.\,1943 Wien@\textsc{Sosnosky, Theodor von} (4.\,1.\,1866 Budapest – 3.\,2.\,1943 Wien), \emph{Schriftsteller}!Neue Erzählungsliteratur@\strich\emph{Neue Erzählungsliteratur}|pwv}{}\ledrightnote{{$\rightarrow$}\emph{\textcolor{green}{Neue Erzählungsliteratur}}} Herrn \textcolor{blue}{Pötzl}\pwindex{Pötzl, Eduard 17.\,3.\,1851 Wien – 20.\,8.\,1914 Mödling@\textsc{Pötzl, Eduard} (17.\,3.\,1851 Wien – 20.\,8.\,1914 Mödling), \emph{Schriftsteller, Journalist}|pw}{}\ledrightnote{\textcolor{blue}{Eduard Pötzl}} zuschicken, der sie umgehend brachte. Ich habe bei
               diesem Anlass auch gleich einige andere Bücher, die ich von Autoren und Verlegern
               bekommen, mit\strikeout{\textcolor{gray}{aufnehm}}\introOben{}geno{\geminationm}\introOben{}en, darunter auch Ihres, wie wohl ich dessen Ablehnung fürchtete, da mir dessen
               Besprechung durch \textcolor{blue}{Bahr}\pwindex{Bahr, Hermann 19.\,7.\,1863 Linz – 15.\,1.\,1934 München@\textsc{Bahr, Hermann} (19.\,7.\,1863 Linz – 15.\,1.\,1934 München), \emph{Schriftsteller, Kritiker}|pw}{}\ledrightnote{\textcolor{blue}{Hermann Bahr}} wahrscheinlich schien.
               Diese Befürchtung wuchs, als ich am Tage nach der Einsendung den \label{K_L04232-5v}\edtext{\textcolor{green}{Artikel}\pwindex{W. Fred 29.\,6.\,1879 Wien – 23.\,10.\,1922 Berlin@\textsc{W. Fred} (29.\,6.\,1879 Wien – 23.\,10.\,1922 Berlin), \emph{Schriftsteller, Journalist}!Stille Bücher@\strich\emph{Stille Bücher}|pwv}{}\ledrightnote{{$\rightarrow$}\emph{\textcolor{green}{Stille Bücher}}}{ }\textcolor{blue}{Freds}\pwindex{W. Fred 29.\,6.\,1879 Wien – 23.\,10.\,1922 Berlin@\textsc{W. Fred} (29.\,6.\,1879 Wien – 23.\,10.\,1922 Berlin), \emph{Schriftsteller, Journalist}|pw}{}\ledrightnote{\textcolor{blue}{W. Fred}}}{\lemma{\textnormal{\emph{Artikel Freds}}}\Cendnote{\textnormal{\textcolor{blue}{W. Fred}\pwindex{W. Fred 29.\,6.\,1879 Wien – 23.\,10.\,1922 Berlin@\textsc{W. Fred} (29.\,6.\,1879 Wien – 23.\,10.\,1922 Berlin), \emph{Schriftsteller, Journalist}|pwk}: \emph{\textcolor{green}{Stille Bücher}\pwindex{W. Fred 29.\,6.\,1879 Wien – 23.\,10.\,1922 Berlin@\textsc{W. Fred} (29.\,6.\,1879 Wien – 23.\,10.\,1922 Berlin), \emph{Schriftsteller, Journalist}!Stille Bücher@\strich\emph{Stille Bücher}|pwk}}. In: \emph{\textcolor{green}{Neues Wiener
                        Tagblatt}\pwindex{Neues Wiener Tagblatt@\emph{Neues Wiener Tagblatt}|pwk}}, Jg. 35, Nr. 164, 17. 6. 1901, S. 10.
                  Neben \emph{\textcolor{green}{Bertha Garlan}\pwindex{Schnitzler, Arthur 15. 5. 1862 Wien – 21. 10. 1931 ebd.@\textsc{Schnitzler, Arthur} (15. 5. 1862 Wien – 21. 10. 1931 ebd.), \emph{Schriftsteller, Mediziner}!Frau Bertha Garlan. Roman@\strich\emph{Frau Bertha Garlan. Roman}|pwk}} wurden die
                  Erzählsammlungen \emph{\textcolor{green}{Die Troika}\pwindex{David, Jakob Julius 6.\,2.\,1859 Hranice – 20.\,11.\,1906 Wien@\textsc{David, Jakob Julius} (6.\,2.\,1859 Hranice – 20.\,11.\,1906 Wien), \emph{Schriftsteller, Journalist}!Troika@\strich\emph{Die Troika}|pwk}} von \textcolor{blue}{Jakob Julius David}\pwindex{David, Jakob Julius 6.\,2.\,1859 Hranice – 20.\,11.\,1906 Wien@\textsc{David, Jakob Julius} (6.\,2.\,1859 Hranice – 20.\,11.\,1906 Wien), \emph{Schriftsteller, Journalist}|pwk} und \emph{\textcolor{green}{Aus Spätherbsttagen}\pwindex{Ebner-Eschenbach, Marie von 13.\,9.\,1830 Zdislavice – 12.\,3.\,1916 Wien@\textsc{Ebner-Eschenbach, Marie von} (13.\,9.\,1830 Zdislavice – 12.\,3.\,1916 Wien), \emph{Schriftstellerin, Schriftstellerin}!Aus Spätherbsttagen. Erzählungen@\strich\emph{Aus Spätherbsttagen. Erzählungen}|pwk}} von \textcolor{blue}{Marie von Ebner-Eschenbach}\pwindex{Ebner-Eschenbach, Marie von 13.\,9.\,1830 Zdislavice – 12.\,3.\,1916 Wien@\textsc{Ebner-Eschenbach, Marie von} (13.\,9.\,1830 Zdislavice – 12.\,3.\,1916 Wien), \emph{Schriftstellerin, Schriftstellerin}|pwk} rezensiert.}}}\label{K_L04232-5} fand; wider Erwarten machte
                  H. \textcolor{blue}{Pötzl}\pwindex{Pötzl, Eduard 17.\,3.\,1851 Wien – 20.\,8.\,1914 Mödling@\textsc{Pötzl, Eduard} (17.\,3.\,1851 Wien – 20.\,8.\,1914 Mödling), \emph{Schriftsteller, Journalist}|pw}{}\ledrightnote{\textcolor{blue}{Eduard Pötzl}} eine Ausnahme und liess aus
               besonderer Rücksicht zwei \textcolor{green}{Besprechungen}\pwindex{W. Fred 29.\,6.\,1879 Wien – 23.\,10.\,1922 Berlin@\textsc{W. Fred} (29.\,6.\,1879 Wien – 23.\,10.\,1922 Berlin), \emph{Schriftsteller, Journalist}!Stille Bücher@\strich\emph{Stille Bücher}|pwv}\pwindex{Sosnosky, Theodor von 4.\,1.\,1866 Budapest – 3.\,2.\,1943 Wien@\textsc{Sosnosky, Theodor von} (4.\,1.\,1866 Budapest – 3.\,2.\,1943 Wien), \emph{Schriftsteller}!Neue Erzählungsliteratur@\strich\emph{Neue Erzählungsliteratur}|pwv}{}\ledrightnote{{$\rightarrow$}\emph{\textcolor{green}{Stille Bücher}}{\newline}{$\rightarrow$}\emph{\textcolor{green}{Neue Erzählungsliteratur}}} desselben \textcolor{green}{Buches}\pwindex{Schnitzler, Arthur 15. 5. 1862 Wien – 21. 10. 1931 ebd.@\textsc{Schnitzler, Arthur} (15. 5. 1862 Wien – 21. 10. 1931 ebd.), \emph{Schriftsteller, Mediziner}!Frau Bertha Garlan. Roman@\strich\emph{Frau Bertha Garlan. Roman}|pwv}{}\ledrightnote{{$\rightarrow$}\emph{\textcolor{green}{Frau Bertha Garlan. Roman}}} zu; nur kürzte er die \textcolor{green}{meinige}\pwindex{Sosnosky, Theodor von 4.\,1.\,1866 Budapest – 3.\,2.\,1943 Wien@\textsc{Sosnosky, Theodor von} (4.\,1.\,1866 Budapest – 3.\,2.\,1943 Wien), \emph{Schriftsteller}!Neue Erzählungsliteratur@\strich\emph{Neue Erzählungsliteratur}|pwv}{}\ledrightnote{{$\rightarrow$}\emph{\textcolor{green}{Neue Erzählungsliteratur}}} sehr begreiflicher Weise. Ich habe
               Ihnen die Genesis dieser \textcolor{green}{Recension}\pwindex{Sosnosky, Theodor von 4.\,1.\,1866 Budapest – 3.\,2.\,1943 Wien@\textsc{Sosnosky, Theodor von} (4.\,1.\,1866 Budapest – 3.\,2.\,1943 Wien), \emph{Schriftsteller}!Neue Erzählungsliteratur@\strich\emph{Neue Erzählungsliteratur}|pwv}{}\ledrightnote{{$\rightarrow$}\emph{\textcolor{green}{Neue Erzählungsliteratur}}} so weitläufig erzählt, auf die Gefahr hin, als langweiliger
               Schwätzer zu erscheinen; aber ich will nicht, das Sie etwa glauben, ich hätte Ihr \textcolor{green}{Buch}\pwindex{Schnitzler, Arthur 15. 5. 1862 Wien – 21. 10. 1931 ebd.@\textsc{Schnitzler, Arthur} (15. 5. 1862 Wien – 21. 10. 1931 ebd.), \emph{Schriftsteller, Mediziner}!Frau Bertha Garlan. Roman@\strich\emph{Frau Bertha Garlan. Roman}|pw}{}\ledrightnote{\textcolor{green}{Frau Bertha Garlan. Roman}} nur so obenhin abgethan, eine Befürchtung,
               die um so näher liegt, als Sie, der Sie mich ja nicht näher kennen, leicht glauben
               könnten, ich \substVorne{}\textsuperscript{\textcolor{gray}{×}\-\textcolor{gray}{×}\-\textcolor{gray}{×}\-\textcolor{gray}{×}\-\textcolor{gray}{×}\-\textcolor{gray}{×}\-\textcolor{gray}{×}\-\textcolor{gray}{×}\-\textcolor{gray}{×}\-\textcolor{gray}{×}\-\textcolor{gray}{×}\-\textcolor{gray}{×}\-\textcolor{gray}{×}\-\textcolor{gray}{×}\-\textcolor{gray}{×}\-\textcolor{gray}{×}\-\textcolor{gray}{×}\-\textcolor{gray}{×}\-\textcolor{gray}{×}\-\textcolor{gray}{×}\-\textcolor{gray}{×}\-\textcolor{gray}{×}{ }\textcolor{gray}{×}\-\textcolor{gray}{×}\-\textcolor{gray}{×}\-\textcolor{gray}{×}\-\textcolor{gray}{×}\-\textcolor{gray}{×}\-\textcolor{gray}{×}\-\textcolor{gray}{×}\-\textcolor{gray}{×}\-\textcolor{gray}{×}\-\textcolor{gray}{×}\-\textcolor{gray}{×}}\substDazwischen{}wäre in meinen Urteile\substHinten{} durch \textcolor{green}{Lt. Gustl}\pwindex{Schnitzler, Arthur 15. 5. 1862 Wien – 21. 10. 1931 ebd.@\textsc{Schnitzler, Arthur} (15. 5. 1862 Wien – 21. 10. 1931 ebd.), \emph{Schriftsteller, Mediziner}!Lieutenant Gustl. Novelle@\strich\emph{Lieutenant Gustl. Novelle}|pw}{}\ledrightnote{\textcolor{green}{Lieutenant Gustl. Novelle}} beeinflusst. Was die
               ausführlichere \textcolor{green}{Besprechung}\pwindex{Sosnosky, Theodor von 4.\,1.\,1866 Budapest – 3.\,2.\,1943 Wien@\textsc{Sosnosky, Theodor von} (4.\,1.\,1866 Budapest – 3.\,2.\,1943 Wien), \emph{Schriftsteller}!Bücher und Zeitschriftenschau [Frau Bertha Garlan]@\strich\emph{Bücher und Zeitschriftenschau [Frau Bertha Garlan]}|pwv}{}\ledrightnote{{$\rightarrow$}\emph{\textcolor{green}{Bücher und Zeitschriftenschau [Frau Bertha Garlan]}}}
               anbelangt, so werden Sie sie vermutlich durch den \textcolor{brown}{Observer}\orgindex{Observer. Alexander Weigl’s Unternehmen für Zeitungssausschnitte@Observer. Alexander Weigl’s Unternehmen für Zeitungssausschnitte|pw}{}\ledrightnote{\textcolor{brown}{Observer. Alexander Weigl’s Unternehmen für Zeitungssausschnitte}} früher \strikeout{wo} zu Gesichte bekommen als
               ich; ich bitte daher freundlichst zu entschuldigen, wenn ich sie Ihnen vielleicht
               sehr lange nicht werde zusenden können. \textcolor{green}{Lt.
                  Gustl}\pwindex{Schnitzler, Arthur 15. 5. 1862 Wien – 21. 10. 1931 ebd.@\textsc{Schnitzler, Arthur} (15. 5. 1862 Wien – 21. 10. 1931 ebd.), \emph{Schriftsteller, Mediziner}!Lieutenant Gustl. Novelle@\strich\emph{Lieutenant Gustl. Novelle}|pw}{}\ledrightnote{\textcolor{green}{Lieutenant Gustl. Novelle}} habe ich ebenfalls besprochen und zwar in einem eigenen \textcolor{green}{Artikel}\pwindex{Sosnosky, Theodor von 4.\,1.\,1866 Budapest – 3.\,2.\,1943 Wien@\textsc{Sosnosky, Theodor von} (4.\,1.\,1866 Budapest – 3.\,2.\,1943 Wien), \emph{Schriftsteller}!Fall Schnitzler. Eine unbefangene Betrachtung@\strich\emph{Der Fall Schnitzler. Eine unbefangene Betrachtung}|pwv}{}\ledrightnote{{$\rightarrow$}\emph{\textcolor{green}{Der Fall Schnitzler. Eine unbefangene Betrachtung}}}, wobei ich bemüht war, den grossen
               künstlerischen Vorzügen der \textcolor{green}{Erzählung}\pwindex{Schnitzler, Arthur 15. 5. 1862 Wien – 21. 10. 1931 ebd.@\textsc{Schnitzler, Arthur} (15. 5. 1862 Wien – 21. 10. 1931 ebd.), \emph{Schriftsteller, Mediziner}!Lieutenant Gustl. Novelle@\strich\emph{Lieutenant Gustl. Novelle}|pwv}{}\ledrightnote{{$\rightarrow$}\emph{\textcolor{green}{Lieutenant Gustl. Novelle}}} streng gerecht zu we\introOben{}r\introOben{}den: mit meiner
                  persön{\pb}lichen Empfindung bei der Lectüre habe ich aber keineswegs
               hinter dem Berge gehalten. Ob der \textcolor{green}{Artikel}\pwindex{Sosnosky, Theodor von 4.\,1.\,1866 Budapest – 3.\,2.\,1943 Wien@\textsc{Sosnosky, Theodor von} (4.\,1.\,1866 Budapest – 3.\,2.\,1943 Wien), \emph{Schriftsteller}!Fall Schnitzler. Eine unbefangene Betrachtung@\strich\emph{Der Fall Schnitzler. Eine unbefangene Betrachtung}|pwv}{}\ledrightnote{{$\rightarrow$}\emph{\textcolor{green}{Der Fall Schnitzler. Eine unbefangene Betrachtung}}} aber auch \label{K_L04232-6v}\edtext{erscheinen}{\lemma{\textnormal{\emph{erscheinen}}}\Cendnote{\textnormal{\textcolor{blue}{Theodor von Sosnosky}\pwindex{Sosnosky, Theodor von 4.\,1.\,1866 Budapest – 3.\,2.\,1943 Wien@\textsc{Sosnosky, Theodor von} (4.\,1.\,1866 Budapest – 3.\,2.\,1943 Wien), \emph{Schriftsteller}|pwk}: \emph{\textcolor{green}{Der Fall Schnitzler. Eine unbefangene Betrachtung}\pwindex{Sosnosky, Theodor von 4.\,1.\,1866 Budapest – 3.\,2.\,1943 Wien@\textsc{Sosnosky, Theodor von} (4.\,1.\,1866 Budapest – 3.\,2.\,1943 Wien), \emph{Schriftsteller}!Fall Schnitzler. Eine unbefangene Betrachtung@\strich\emph{Der Fall Schnitzler. Eine unbefangene Betrachtung}|pwk}}. In:
                        \emph{\textcolor{green}{Neue Bahnen}\pwindex{Neue Bahnen. Zeitschrift für Kunst und öffentliches Leben@\emph{Neue Bahnen. Zeitschrift für Kunst und öffentliches Leben}|pwk}}, Jg. 1, Nr. 18,
                        1. 10. 1901, S. 508–514.}}}\label{K_L04232-6} wird, weiss ich nicht;
               für \textcolor{pink}{Wien}\oindex{Wien@\textbf{Wien}, \emph{Verwaltungsgebiet}|pw}{}\ledrightnote{\textcolor{pink}{Wien}} hat er wenig Chancen, denn der
               überwiegende Teil der Presse steht natürlich principiell auf \substVorne{}\textsuperscript{ihr}\substDazwischen{}Ihr\substHinten{}er Seite und würde meiner Auffassung sicher nicht Raum gewähren, der andere
               Teil aber würde \introOben{}–\introOben{} ebenso principiell \introOben{}–\introOben{} Ihr \uline{Gegner} \introOben{}–\introOben{}
               meine \textcolor{green}{Arbeit}\pwindex{Sosnosky, Theodor von 4.\,1.\,1866 Budapest – 3.\,2.\,1943 Wien@\textsc{Sosnosky, Theodor von} (4.\,1.\,1866 Budapest – 3.\,2.\,1943 Wien), \emph{Schriftsteller}!Fall Schnitzler. Eine unbefangene Betrachtung@\strich\emph{Der Fall Schnitzler. Eine unbefangene Betrachtung}|pwv}{}\ledrightnote{{$\rightarrow$}\emph{\textcolor{green}{Der Fall Schnitzler. Eine unbefangene Betrachtung}}} zweifellos zu
               freundlich finden und \uline{darum} ablehnen.\pend
           
\pstart
           Eingesendet ist die \textcolor{green}{Arbeit}\pwindex{Sosnosky, Theodor von 4.\,1.\,1866 Budapest – 3.\,2.\,1943 Wien@\textsc{Sosnosky, Theodor von} (4.\,1.\,1866 Budapest – 3.\,2.\,1943 Wien), \emph{Schriftsteller}!Fall Schnitzler. Eine unbefangene Betrachtung@\strich\emph{Der Fall Schnitzler. Eine unbefangene Betrachtung}|pwv}{}\ledrightnote{{$\rightarrow$}\emph{\textcolor{green}{Der Fall Schnitzler. Eine unbefangene Betrachtung}}}
               übrigens.\pend
           
\pstart
           Mit den verbindlichsten Empfehlungen{\\[\baselineskip]}hochachtungsvoll{\\[\baselineskip]}\spacefill\mbox{Theodor vSosnosky}\pend
           \leftskip=0em{}\selectlanguage{ngerman}\endnumbering\briefempfaengerindex{Schnitzler, Arthur@\textsc{Schnitzler, Arthur}!zzzSosnosky, Theodor von@\emph{von Theodor von Sosnosky}!1901-06-212@{21. 6. 1901}|)be}\mylabel{L04232h}
\begin{anhang}
\end{anhang}\newcommand{\dateiname}{L04232}\newcommand{\titel}{Theodor von Sosnosky an Arthur Schnitzler, 21. 6. 1901}\newcommand{\editorInnen}{Herausgegeben von Jahnke, SelmaMüller, Martin Anton}\input{../tex-inputs/latex-korrekturansicht-abspann}
